\section{A decision-theoretic value of acquired oblique geometry}\label{sec:decision99}
An approximation separation matters economically only when it changes an implemented decision. We give such a realization for the original oblique continuation, without altering the retained investment comparison. This is an exactly solved economic mechanism, not an empirical calibration or a timing experiment.

A decision maker observes $x=(x_1,x_2)\in[0,1]^2$ and a posted price $p\in[0,3]$. Purchasing one durable claim yields next-date benefit
\begin{equation}\label{eq:benefit99}
 v(x)=1+(x_1+x_2-1)_+.
\end{equation}
The feasible action is $b\in\{0,1\}$, and total cost is $b\{p-v(x)\}$. This is a two-date Bellman problem with a deterministic terminal payoff and unit discount. A stored continuation $h(x)\in[0,3]$ implements $b_h(x,p)=\boldsymbol1\{p\le h(x)\}$. The original optimum buys at $p\le v(x)$. A price tie has zero loss. Let states be uniform on the square and prices uniform on $[0,3]$, independently.

\begin{theorem}[Exact policy-loss separation]\label{thm:decision99}
For every measurable $h:[0,1]^2\to[0,3]$, the conditional and unconditional policy regrets are
\begin{equation}\label{eq:decisionloss99}
 \mathbb E_p[J^{b_h}(x,p)-J^*(x,p)]
       =\frac{(h(x)-v(x))^2}{6},\qquad
 \mathbb E[J^{b_h}-J^*]=\frac{\|h-v\|_{L^2}^2}{6}.
\end{equation}
Among arbitrary square-integrable additive continuations and among polynomials of total degree at most two, respectively, the optimal expected regrets are
\begin{equation}\label{eq:separation99}
 \frac1{540}\quad\text{and}\quad\frac1{5400}.
\end{equation}
They are attained by the admissible functions
\begin{align}
 h_A(x)&=\frac56+\frac{x_1^2+x_2^2}{2},\label{eq:additive99}\\
 h_Q(x)&=\frac{17}{15}-\frac35(x_1+x_2)
              +\frac12(x_1^2+x_2^2)+\frac65x_1x_2.\label{eq:quadratic99}
\end{align}
A single oblique ReLU with the affine constant in \eqref{eq:benefit99} attains zero regret. A conventional implementation of the identical oblique hinge also attains zero regret and the same decisions.
\end{theorem}

The shift by one is part of the economic construction. It keeps the projection optima inside the posted-price support, so clipping does not secretly enlarge a comparator class. The proof computes the additive conditional-expectation projection and the six-dimensional polynomial Gram system exactly. It establishes a policy comparison against the best member of each specified class, rather than comparing arbitrarily chosen fits. The scalar price integration is the step that turns continuation error into action loss; no such identity is presumed in the investment economy.

The sup-norm comparison also has a decision interpretation. Under the same conditional price law, the smallest worst-state expected regrets for the two classes are $1/96$ and $1/1536$, respectively. These follow from the sharp errors $1/4$ and $1/16$ in Proposition~\ref{prop:ridge97}, with its explicit feasible attaining functions shifted by one. The population optima in \eqref{eq:additive99}--\eqref{eq:quadratic99} are different from the minimax optima. An acquisition objective must specify which context criterion it is intended to serve.

\subsection{Acquisition, conventional geometry and purchase}
The mechanism identifies a value of the oblique continuation geometry, not a value of relabeling its formula as a network. A piecewise-affine conventional representation on the two triangles separated by $x_1+x_2=1$ is also exact. Convex regression with that hinge already in its dictionary contains the same exact continuation. Those controls are included in the exact comparison. Discovering the separating hyperplane from unknown data, fitting a larger dictionary or constructing an adaptive triangulation are different acquisition problems; their costs cannot be inferred from this formula-access calculation.

Suppose the exact oblique service has additional fixed acquisition cost $K\ge0$, incremental per-decision computation cost $c$, and fee $F\ge0$, all expressed in the units of this model's cost. Against the best quadratic continuation, the exact gain from $m$ independent decisions is
\begin{equation}\label{eq:exactpurchase99}
 G_Q(m)=m(1/5400-c)-K-F.
\end{equation}
When $c<1/5400$, the first strictly beneficial workload is
$\lfloor(K+F)/(1/5400-c)\rfloor+1$. This is a nonzero, exactly identified policy benefit with an explicit acquisition threshold. It is not an estimate of the price of computation in the retained investment application. Equal hinge and neural implementations share this gain when their complete costs agree.

\begin{table}[htbp]
\caption{Exact original-law policy loss in the oblique-continuation mechanism}\label{tab:decision99}
\centering\small
\begin{tabular}{@{}p{.52\linewidth}p{.17\linewidth}p{.22\linewidth}@{}}
\toprule
Continuation class or supplied geometry & Optimal mean regret & Information boundary\\\midrule
Arbitrary square-integrable additive functions & $1/540$ & Population projection\\
Total-degree-two polynomials & $1/5400$ & Six exact moments\\
One oblique ReLU and affine constant & $0$ & Supplied hinge\\
Identical conventional oblique hinge & $0$ & Same supplied hinge\\
Piecewise affine on the two separating triangles & $0$ & Same separating line\\
Convex dictionary containing the oblique hinge & $0$ & Contains exact function\\\bottomrule
\end{tabular}
\par\smallskip\parbox{.96\linewidth}{\footnotesize Notes: States are uniform on $[0,1]^2$ and independent prices uniform on $[0,3]$. Loss is evaluated against the original purchase optimum. The last four rows are equivalent exact geometric realizations, not four independent fitted experiments. The quadratic row uses six right-side moments and its known Gram matrix; moment acquisition is not free when the continuation formula is unknown. No timing, empirical calibration or dimension-wide neural dominance is inferred.}
\end{table}

The table is generated from two separately implemented exact-arithmetic calculations. The production path assembles monomial moments and solves normal equations by rational elimination. The independent path integrates polynomials after a triangular change of variables and uses an orthogonal shifted-Legendre basis. It imports neither the production integrator nor its solver. Both reproduce the fitted function, the original-law loss and the range checks. This independently checks acquisition for the stated finite-dimensional mechanism; it does not independently refit the R97 nonlinear dictionary search or complete the separately frozen R98 campaign.
